\begin{proposition}[6.10.3]
\label{IV.2.s6.10.3-fr}
Con las notaciones de \textup{(6.10.1)}, supongamos que $\mathcal F$ es
normalmente plano a lo largo de $Y$. Entonces:
\begin{enumerate}
\item[\textup{(i)}] $Y\cap\operatorname{Supp}(\mathcal F)$ es una parte a la vez
abierta y cerrada de $Y$ \textup{(es decir, una} unión de componentes conexas
de $Y$).
\item[\textup{(ii)}] Si $\mathcal J$ es localmente nilpotente,
$\operatorname{Supp}(\mathcal F)$ es una parte abierta y cerrada de $X$, y
para todo $x\in\operatorname{Supp}(\mathcal F)$ se tiene
\end{enumerate}
\begin{align}
\dim(\mathcal F_x)&=\dim(\mathcal O_{Y,x})
  \tag{6.10.3.1}\label{IV.2.s6.10.3.1-fr}\\
\operatorname{prof}(\mathcal F_x)&=\operatorname{prof}(\mathcal O_{Y,x})
  \tag{6.10.3.2}\label{IV.2.s6.10.3.2-fr}\\
\operatorname{coprof}(\mathcal F_x)&=\operatorname{coprof}(\mathcal O_{Y,x}).
  \tag{6.10.3.3}\label{IV.2.s6.10.3.3-fr}
\end{align}
\end{proposition}

\begin{proof}
\begin{enumerate}
\item[(i)] Con las notaciones de (6.10.1), se tiene
$Y\cap\operatorname{Supp}(\mathcal F)=\operatorname{Supp}(\mathcal G_0)$
(\textbf I, 9.1.13), y como, por hipótesis, $\mathcal G_0$ es un
$\mathcal O_Y$-Módulo localmente libre y finitamente generado, su soporte es a la vez
abierto y cerrado en $Y$.

\oldpage[IV-2]{156}

\item[(ii)] La hipótesis de que $\mathcal J$ es localmente nilpotente implica que
los espacios subyacentes de $X$ y de $Y$ coinciden; de ahí se siguen las dos primeras
afirmaciones de (ii), teniendo en cuenta (5.1.12.1). Queda por demostrar (6.10.3.2),
pues de ella se deduce inmediatamente (6.10.3.3) por diferencia. Sea
$(f_i)_{1\leq i\leq p}$ una sucesión finita de elementos del ideal maximal de
$\mathcal O_{X,x}$ cuyas imágenes en
$\mathcal O_{Y,x}=\mathcal O_{X,x}/\mathcal J_x$ forman una sucesión
$\mathcal O_{Y,x}$-regular \emph{maximal}; la hipótesis sobre $\mathcal F$ y
$\mathcal J$ implica que el $\mathcal O_{Y,x}$-módulo
$\operatorname{gr}_{\mathcal J_x}^{\bullet}(\mathcal F_x)$ es libre y finitamente
generado; por tanto, la sucesión $(f_i)$ es
$\operatorname{gr}_{\mathcal J_x}^{\bullet}(\mathcal F_x)$-regular. Se
deduce que esta sucesión es también $\mathcal F_x$-\emph{regular} (\textbf 0, 15.1.19). Sea,
por otra parte, $n$ el mayor entero tal que
$\mathcal F_x^{(n)}=\mathcal J_x^n\mathcal F_x\ne0$. La hipótesis implica
también que
$\operatorname{gr}_{\mathcal J_x}^{\bullet}
(\mathcal F_x/\mathcal F_x^{(n)})$ es libre y finitamente generado, por lo que la sucesión
$(f_i)$ es también $(\mathcal F_x/\mathcal F_x^{(n)})$-\emph{regular}
(\emph{loc. cit.}). Aplicando el lema (3.4.1.4), se obtiene por inducción
sobre $i$ una sucesión exacta
\[
0\longrightarrow
\mathcal F_x^{(n)}/\left(\sum_{i=1}^{p}f_i\mathcal F_x^{(n)}\right)
\longrightarrow
\mathcal F_x/\left(\sum_{i=1}^{p}f_i\mathcal F_x\right).
\]
Pero la hipótesis implica que $\mathcal F_x^{(n)}$ es un
$\mathcal O_{Y,x}$-módulo libre, finitamente generado y $\ne0$; por tanto,
$\mathcal F_x^{(n)}/(\sum_{i=1}^{p}f_i\mathcal F_x^{(n)})$ es isomorfo a
un módulo de la forma
$(\mathcal O_{Y,x}/(\sum_{i=1}^{p}f_i\mathcal O_{Y,x}))^m$ con $m>0$; como
$\operatorname{prof}(\mathcal O_{Y,x}/(\sum_{i=1}^{p}f_i\mathcal O_{Y,x}))=0$
(\textbf 0, 16.4.6) se tiene también, por la caracterización (\textbf 0, 16.4.6, (i)) de los módulos
de profundidad nula,
$\operatorname{prof}(\mathcal F_x^{(n)}/(\sum_{i=1}^{p}f_i\mathcal F_x^{(n)}))=0$,
y después
$\operatorname{prof}(\mathcal F_x/(\sum_{i=1}^{p}f_i\mathcal F_x))=0$;
como los $(f_i)_x$ pertenecen al ideal maximal de $\mathcal O_{Y,x}$ y forman
una sucesión $\mathcal F_x$-\emph{regular}, esto demuestra que se tiene
$\operatorname{prof}(\mathcal F_x)=p$ (\textbf 0, 16.4.6, (ii)).
\end{enumerate}
\end{proof}

\begin{corollary}[6.10.4]
\label{IV.2.s6.10.4-fr}
Sea $U_{S_n}(\mathcal F)$ \textup{(resp.} $U_{C_n}(\mathcal F)$) el conjunto de los
$x\in X$ tales que $\mathcal F$ satisface $(S_n)$ en el punto $x$ \textup{(resp.} tales que
$\operatorname{coprof}(\mathcal F_x)\leq n$). Si $\mathcal F$ es normalmente
plano a lo largo de $Y$, si $\operatorname{Supp}(\mathcal F)=X$ y si $\mathcal J$
es localmente nilpotente, se tiene
$U_{S_n}(\mathcal F)=U_{S_n}(\mathcal O_Y)$ y
$U_{C_n}(\mathcal F)=U_{C_n}(\mathcal O_Y)$.
\end{corollary}
